\documentclass[10pt,a4paper]{amsart}
\usepackage[margin=19mm]{geometry}
\usepackage{amsmath,amssymb,mathtools}
\usepackage[scr=rsfs,cal=euler]{mathalfa}
\usepackage{xfrac}
\usepackage[unicode,hidelinks,bookmarks=false]{hyperref}
\newcommand{\ie}{i.e.\ }
\newcommand{\cf}{cf.\ }
\newcommand{\resp}{respectively\ }
\newcommand{\Subsection}{\subsection}
\providecommand{\nobreakdash}{\nobreak\hbox{-}}
\setlength{\emergencystretch}{2em}
\multlinegap0pt
\usepackage{bm}
\usepackage{booktabs}
\usepackage{array}
\usepackage{pgfplots}
\usepackage{pgfplotstable}
\pgfplotsset{compat=1.17}
\setlength{\emergencystretch}{2em}
\numberwithin{equation}{section}
\newtheorem{theorem}{Theorem}[section]
\newtheorem{assumption}{Assumption}[section]
\newtheorem{remark}{Remark}[section]
\newcommand{\ii}{\mathrm{i}}
\newcommand{\dd}{\,\mathrm{d}}
\newcommand{\RR}{\mathbb{R}}
\begin{document}
\title[A High-Order Nystr"om Method for Coupled Boundary Integral E]{A High-Order Nystr\"om Method for Coupled Boundary Integral Equations in Oblique-Incidence Scattering by Impedance Cylinders}
\author{Haochen Liu Qinghao Yu; Haochen Liu; Qinghao Yu}
\date{}
\maketitle
\noindent\emph{Source and licence.} A High-Order Nystr\"om Method for Coupled Boundary Integral Equations in Oblique-Incidence Scattering by Impedance Cylinders. Version arXiv:2606.15616v1 (CC BY 4.0). This material is distributed under the Creative Commons Attribution 4.0 International licence, http://creativecommons.org/licenses/by/4.0/, as recorded in the arXiv OAI licence metadata for this exact version and shown on the abstract page https://arxiv.org/abs/2606.15616v1. Full rights notice: licenses/arxiv-2606.15616-CC-BY-4.0.txt.\par\smallskip
\noindent\textit{Editorial scope.} Counted unit: the original \texttt{theorem} environment \texttt{-} at source lines 143--149 together with its complete proof at lines 151--155; the delivered document keeps source lines 70--156 so that every referenced label and every citation resolves inside it. Native editable source is the paper's own concatenated TeX; the AMS wrapper is editorial formatting only. No publisher class, upstream script or unrelated artwork is redistributed.
\section{Mathematical Model and Boundary Integral Formulation}

\subsection{Axial-component scattering model}

Let $D\subset\RR^2$ be the bounded cross section of an infinitely long cylinder, with smooth boundary $\Gamma=\partial D$, and let $\Omega=\RR^2\setminus\overline D$ be the exterior domain. We use the time factor $\exp(-\ii\omega t)$. The free-space wavenumber is $k>0$, the incidence angle relative to the cylinder axis is $\alpha\in(0,\pi/2]$, and
\begin{equation}
\beta=k\cos\alpha,\qquad \kappa=k\sin\alpha
\end{equation}
denote the axial propagation constant and transverse wavenumber, respectively.

After separation of the factor $\exp(\ii\beta z)$, the scattered axial electric and magnetic components
\[
u=E_z^s,\qquad v=H_z^s
\]
satisfy
\begin{equation}
\Delta u+\kappa^2u=0,\qquad \Delta v+\kappa^2v=0\qquad \text{in }\Omega,
\label{eq:helmholtz}
\end{equation}
together with the Sommerfeld radiation condition for each component. Let $\nu$ be the exterior unit normal, let $\tau$ be the positively oriented unit tangent, and let $\partial_s=\tau\cdot\nabla$ denote arclength differentiation. With a fixed orientation convention, the axial-component form of the normalized Leontovich condition is written as
\begin{align}
\partial_\nu u+\ii\eta u-\mu\,\partial_s v&=f_1 \quad \text{on }\Gamma, \label{eq:bc1}\\
\partial_\nu v+\ii\eta v+\mu\,\partial_s u&=f_2 \quad \text{on }\Gamma. \label{eq:bc2}
\end{align}
Here $\eta$ is the normalized surface impedance, which may be a smooth complex-valued function on $\Gamma$, and $\mu$ is an oblique-incidence coupling coefficient determined by the normalization and the material parameters. The signs of the off-diagonal terms depend on the orientation of $\tau$; reversing the orientation changes both signs and leaves the coupled structure unchanged. We assume a passive impedance in the sense appropriate to the chosen time convention, and in the experiments use $\operatorname{Re}\eta\ge0$ and $\operatorname{Im}\eta\ge0$.

\subsection{Layer potentials and the coupled system}

Let
\begin{equation}
\Phi_\kappa(x,y)=\frac{\ii}{4}H_0^{(1)}(\kappa |x-y|)
\end{equation}
be the outgoing fundamental solution of the two-dimensional Helmholtz equation. For a boundary density $\varphi$ define the single-layer potential
\begin{equation}
(\mathcal S\varphi)(x)=\int_\Gamma \Phi_\kappa(x,y)\varphi(y)\,\dd s_y,\qquad x\in\Omega.
\end{equation}
On $\Gamma$ we write
\begin{align}
(S\varphi)(x)&=\int_\Gamma \Phi_\kappa(x,y)\varphi(y)\,\dd s_y,\\
(K'\varphi)(x)&=\operatorname{p.v.}\int_\Gamma \partial_{\nu_x}\Phi_\kappa(x,y)\varphi(y)\,\dd s_y,\\
(T\varphi)(x)&=\partial_s(S\varphi)(x).
\end{align}
The operator $S$ maps $H^{-1/2}(\Gamma)$ to $H^{1/2}(\Gamma)$, while $K'$ and $T$ are naturally interpreted on Sobolev trace spaces of order $-1/2$ after the standard embeddings and principal-value interpretation. With our convention for the exterior normal trace of the single-layer potential,
\begin{equation}
\partial_\nu^+\mathcal S\varphi=\left(\frac12 I+K'\right)\varphi .
\end{equation}

We represent the scattered fields by two single-layer potentials,
\begin{equation}
u=\mathcal S\varphi_1,\qquad v=\mathcal S\varphi_2.
\end{equation}
Substitution into \eqref{eq:bc1}--\eqref{eq:bc2} gives the coupled boundary integral equation
\begin{equation}
\mathcal A
\begin{bmatrix}\varphi_1\\ \varphi_2\end{bmatrix}
=
\begin{bmatrix}f_1\\ f_2\end{bmatrix},
\qquad
\mathcal A=
\begin{bmatrix}
L&-\mu T\\
\mu T&L
\end{bmatrix},
\label{eq:bie}
\end{equation}
where
\begin{equation}
L=\frac12 I+K'+\ii M_\eta S .
\end{equation}
Here $M_\eta$ is multiplication by $\eta$, so that $M_\eta S\varphi=\eta(\cdot)(S\varphi)(\cdot)$. The diagonal blocks are scalar impedance boundary integral operators. The off-diagonal blocks contain the tangential derivative operator generated by the oblique-incidence coupling.

\subsection{Fredholm framework}


\begin{theorem}[Fredholm solvability framework]
Let $\Gamma$ be a sufficiently smooth closed curve and let $\eta$ be a sufficiently smooth passive impedance. Assume that the single-layer representation is not affected by an internal non-resonance obstruction and that the principal symbol of the coupled boundary operator is elliptic in the relevant Sobolev product space. Then
\[
\mathcal A:H^{-1/2}(\Gamma)^2\to H^{-1/2}(\Gamma)^2
\]
is Fredholm of index zero. If the corresponding exterior impedance scattering problem is uniquely solvable, then the boundary integral equation \eqref{eq:bie} is uniquely solvable.
\end{theorem}

\begin{proof}
We give the standard proof outline, since a complete treatment requires the full boundary pseudodifferential calculus. Consider first the homogeneous equation $\mathcal A\varphi=0$ and let $u=\mathcal S\varphi_1$, $v=\mathcal S\varphi_2$. The jump relations imply that the exterior traces satisfy the homogeneous version of \eqref{eq:bc1}--\eqref{eq:bc2}. Thus $(u,v)$ solves the exterior coupled Helmholtz problem with the radiation condition. Applying Green's identity on $\Omega\cap B_R$ and letting $R\to\infty$, the radiation condition and the passivity of the impedance yield vanishing radiated energy. Rellich's lemma then implies that the far-field pattern is zero, and unique continuation gives $u=v=0$ in the exterior.

It remains to connect the vanishing exterior fields to the densities. The boundary traces and the single-layer jump relations show that the densities solve an associated interior homogeneous problem. Under the stated internal non-resonance condition this implies $\varphi_1=\varphi_2=0$. Hence the nullspace is trivial. Fredholmness follows from the decomposition of $\mathcal A$ into the scalar impedance part, smoothing contributions of the single-layer trace, and the principal-value tangential term, whose coupled principal symbol is assumed elliptic. Since the index is zero, the Fredholm alternative gives unique solvability.
\end{proof}


\end{document}
